
\documentclass{article}
\usepackage{colortbl}
\usepackage{makecell}
\usepackage{multirow}
\usepackage{supertabular}

\begin{document}

\newcounter{utterance}

\centering \large Interaction Transcript for game `cladder', experiment `full\_v1.5\_default', episode 8235 with BSU{-}SLIM\_\_Qwen3.5{-}9B{-}Base\_\_playpen{-}prm{-}9b{-}best\_of\_n.
\vspace{24pt}

{ \footnotesize  \setcounter{utterance}{1}
\setlength{\tabcolsep}{0pt}
\begin{supertabular}{c@{$\;$}|p{.15\linewidth}@{}p{.15\linewidth}p{.15\linewidth}p{.15\linewidth}p{.15\linewidth}p{.15\linewidth}}
    \# & \multicolumn{2}{c}{Player} && \multicolumn{2}{c}{Game Master} \\
    \hline

    \theutterance \stepcounter{utterance}  
    & & & \multicolumn{4}{p{0.6\linewidth}}{
        \cellcolor[rgb]{0.9,0.9,0.9}{
            \makecell[{{p{\linewidth}}}]{
                \texttt{\tiny{[P1$\langle$GM]}}
                \texttt{You are an expert in causal inference. The following question is not a typical commonsense query, but rather a meticulously designed question created by a professor specializing in causal inference, intended to assess the students' mastery of the course content.} \\
\texttt{Imagine a self{-}contained, hypothetical world with only the following conditions, and without any unmentioned factors or causal relationships: Yomx has a direct effect on gwet. Gwet has a direct effect on xevu. For those who are not yomx, the probability of xevu is 68\%. For those who are yomx, the probability of xevu is 40\%. For those who are yomx, would it be less likely to see xevu if the individual was not yomx?} \\
\texttt{Start your answer with 'Yes' or 'No', followed by additional reasoning or evidence to support your explanation.} \\
            }
        }
    }
     \\ \\

    \theutterance \stepcounter{utterance}  
    & \multicolumn{4}{p{0.6\linewidth}}{
        \cellcolor[rgb]{0.561,0.737,0.561}{
            \makecell[{{p{\linewidth}}}]{
                \texttt{\tiny{[P1$\rangle$GM]}}
                \texttt{No} \\
\\ 
\texttt{In this hypothetical world, the causal structure is linear: **Yomx → Gwet → Xevu**.} \\
\\ 
\texttt{To answer the question, we must compare the probability of **Xevu** under two specific conditions:} \\
\texttt{1.  **Condition A (Actual):** The individual *is* Yomx. In this case, the probability of Xevu is explicitly given as **40\%**.} \\
\texttt{2.  **Condition B (Counterfactual):** The individual *was* Yomx, but we intervene to set their state to *not* Yomx.} \\
\\ 
\texttt{Here is the logical breakdown based on the provided rules:} \\
\texttt{*   Since Gwet is the direct cause of Xevu, the probability of Xevu depends on the state of Gwet.} \\
\texttt{*   The problem states: "For those who are yomx, the probability of xevu is 40\%." This implies that being Yomx leads to a specific level of Gwet (let's call it \$Gwet\_\{high\}\$), which results in a 40\% chance of Xevu.} \\
\texttt{*   The problem also states: "For those who are **not** yomx, the probability of xevu is 68\%." This implies that being *not* Yomx leads to a different level of Gwet (let's call it \$Gwet\_\{low\}\$), which results in a 68\% chance of Xevu.} \\
\\ 
\texttt{Therefore, if an individual who is currently Yomx were to be **not** Yomx, their probability of seeing Xevu would shift from 40\% to **68\%**.} \\
\\ 
\texttt{Since 68\% is greater than 40\%, it would be **more** likely (not less likely) to see Xevu if the individual was not Yomx. The presence of Yomx acts as a suppressor in this specific causal chain, reducing the likelihood of the outcome Xevu compared to its absence.} \\
            }
        }
    }
    & & \\ \\

    \theutterance \stepcounter{utterance}  
    & & & \multicolumn{2}{p{0.3\linewidth}}{
        \cellcolor[rgb]{0.9,0.9,0.9}{
            \makecell[{{p{\linewidth}}}]{
                \texttt{\tiny{[GM$|$GM]}}
                \texttt{No} \\
            }
        }
    }
    & & \\ \\

    \theutterance \stepcounter{utterance}  
    & & & \multicolumn{2}{p{0.3\linewidth}}{
        \cellcolor[rgb]{0.9,0.9,0.9}{
            \makecell[{{p{\linewidth}}}]{
                \texttt{\tiny{[GM$|$GM]}}
                \texttt{no} \\
            }
        }
    }
    & & \\ \\

    \theutterance \stepcounter{utterance}  
    & & & \multicolumn{2}{p{0.3\linewidth}}{
        \cellcolor[rgb]{0.9,0.9,0.9}{
            \makecell[{{p{\linewidth}}}]{
                \texttt{\tiny{[GM$|$GM]}}
                \texttt{game\_result = WIN} \\
            }
        }
    }
    & & \\ \\

\end{supertabular}
}

\end{document}
